\begin{figure}[htbp]
\centering
\begin{tikzpicture}[x=1cm,y=1cm,dissertation block/.append style={font=\small\setstretch{1}}]
\node[dissertation block,text width=4.20cm] (lim) at (0,5.5) {Обмеження відомих підходів};
\node[dissertation block,text width=9.00cm] (model) at (0,3.5) {Модель мультимодального представлення\\прихованої загрози};
\node[dissertation block,text width=5.90cm] (svs) at (-3.5,0.4) {Метод аналізу спектральних,\\піксельних і структурних ознак};
\node[dissertation block,text width=5.90cm] (beh) at (3.5,0.4) {Метод поведінкового\\узгодження ознак\\різного походження};
\node[dissertation block,text width=9.00cm] (tech) at (0,-2.8) {Інформаційна технологія виявлення ознак аномального вмісту\\в мультимедійних даних вебсередовища};
\node[dissertation block,text width=5.20cm] (exp) at (0,-5.1) {Експериментальне оцінювання\\та визначення доказових меж};
\draw[dissertation arrow] (lim) -- (model);
\draw[dissertation arrow] (model) -- (svs);
\draw[dissertation arrow] (model) -- (beh);
\draw[dissertation arrow] (svs) -- (tech);
\draw[dissertation arrow] (beh) -- (tech);
\draw[dissertation arrow] (tech) -- (exp);
\end{tikzpicture}
\caption{Логіка переходу від установлених обмежень до експериментального оцінювання}
\label{fig:research_logic}
\end{figure}
